
\documentclass{article}
\usepackage{colortbl}
\usepackage{makecell}
\usepackage{multirow}
\usepackage{supertabular}

\begin{document}

\newcounter{utterance}

\centering \large Interaction Transcript for game `dond', experiment `semi\_en', episode 18 with reflexion{-}Qwen3.5{-}4B.
\vspace{24pt}

{ \footnotesize  \setcounter{utterance}{1}
\setlength{\tabcolsep}{0pt}
\begin{supertabular}{c@{$\;$}|p{.15\linewidth}@{}p{.15\linewidth}p{.15\linewidth}p{.15\linewidth}p{.15\linewidth}p{.15\linewidth}}
    \# & $\;$A & \multicolumn{4}{c}{Game Master} & $\;\:$B\\
    \hline

    \theutterance \stepcounter{utterance}  
    & & \multicolumn{4}{p{0.6\linewidth}}{
        \cellcolor[rgb]{0.9,0.9,0.9}{
            \makecell[{{p{\linewidth}}}]{
                \texttt{\tiny{[P1$\langle$GM]}}
                \texttt{You have attempted this same task before and did not succeed. These are the plans you wrote after those attempts.} \\
\texttt{Trial \#1: New plan: I failed because I miscounted the available items. The set contains only 2 boxes and 1 lemon, but I incorrectly assumed there were 2 grapes. Consequently, I proposed taking 2 grapes, which made my proposal non{-}complementary to the other player's potential claim (since they would need the lemon, and I claimed items they didn't have or in wrong quantities), leading to a zero{-}score outcome. Next time, I must strictly verify the item counts in the prompt before formulating a proposal. I will calculate the total value of all items (2 boxes * 2 + 1 lemon * 4 + 2 grapes * 1 = 10) and propose a split where I take the items with the highest value relative to my total, ensuring the remaining items sum to a positive value for the other player. Specifically, I will propose taking the 2 boxes and 2 grapes (value 6) and give the 1 lemon to the other player (value 4), which is a valid complementary split. I will only submit the proposal in the exact format `{[}Proposal: ...{]}` after confirming the item counts are correct.} \\
\\ 
\texttt{Trial \#2: New plan: I failed because I included "2 grapes" in my proposal, but the available set only contains 1 grape. This made my proposal non{-}complementary to the other player's valid proposal (which took the lemon and 2 boxes, leaving 2 grapes), resulting in a zero score. Next time, I will strictly verify the item counts from the prompt before calculating my proposal. I will calculate the total value (10) and propose taking the 2 boxes and the 1 lemon (value 6), leaving the 2 grapes for the other player (value 2). This creates a complementary split. Crucially, I will ensure that any message sent to the other player (before the final proposal) does not use square brackets, as that format triggers an immediate secret proposal and causes an abort if the proposal is invalid. I will only use the format `{[}Proposal: ...{]}` for the final submission.} \\
\\ 
\texttt{Trial \#3: New plan: I failed because I miscounted the available items, specifically assuming there were 2 grapes when the set only contains 1. This led me to propose taking 2 grapes, which made my proposal non{-}complementary to the other player's valid claim (lemon + 2 boxes), resulting in a zero score. Next time, I will strictly verify the item counts from the prompt before formulating any strategy. I will calculate the total value of all items (2 boxes * 2 + 1 lemon * 4 + 2 grapes * 1 = 10) and propose a split where I take the highest value items available to me: the 1 lemon and 2 boxes (total value 6), leaving the 2 grapes for the other player (total value 2). This ensures a complementary split. Crucially, I will ensure that any message sent to the other player before the final step does not use square brackets, as that format triggers an immediate secret proposal and causes an abort if the proposal is invalid. I will only use the format `{[}Proposal: ...{]}` for the final submission after confirming the item counts are correct.} \\
\\ 
\\ 
\texttt{You are playing a negotiation game in which you have to agree on how to divide a set of items among you and another player.} \\
\\ 
\texttt{Rules:} \\
\texttt{(a) You and the other player are given a set of items. Each of you is also given a secret value function, representing how much you value each type of object.} \\
\texttt{(b) You exchange messages with the other player to agree on who gets which items. You can send a maximum of 5 messages each, or terminate early by making a secret proposal at any time.} \\
\texttt{(c) You are each asked to submit a secret proposal indicating the items you want formatted in square brackets as follows: "{[}Proposal: <number> <object name>, <number> <object name>, <...>{]}"} \\
\texttt{(d) If your proposals are complementary, i.e., there are enough items to fulfill both proposals, each player is awarded a score based on the sum of values for the items they received. Otherwise, both of you get zero points.} \\
\texttt{(e) Your goal is to maximize the score you receive.} \\
\\ 
\texttt{Let us start.} \\
\\ 
\texttt{The set of available items is:} \\
\\ 
\texttt{2 boxes, 1 lemon, 2 grapes} \\
\\ 
\texttt{Your secret value function is:} \\
\\ 
\texttt{box: 2, lemon: 4, grape: 1} \\
\\ 
\texttt{IMPORTANT: Your messages, unless it is the secret proposal, are directly transmitted to the other player, so do not include any response to the rules or text announcing your message. To make a secret proposal, use the indicated format. Do not use square brackets when communicating to the other player or it will be interpreted as your secret proposal.} \\
            }
        }
    }
    & \\ \\

    \theutterance \stepcounter{utterance}  
    & \multicolumn{4}{p{0.6\linewidth}}{
        \cellcolor[rgb]{0.9,0.9,0.9}{
            \makecell[{{p{\linewidth}}}]{
                \texttt{\tiny{[P1$\rangle$GM]}}
                \texttt{Hello! I see the items available: 2 boxes, 1 lemon, and 2 grapes. My values are: box=2, lemon=4, grape=1.} \\
\\ 
\texttt{To maximize my score, I should aim to get the lemon (value 4) and the boxes (value 2 each, total 4), giving me a total of 6. This leaves the 2 grapes (value 2) for the other player. Since 6 + 2 = 8, and the total value of all items is 10, this is a fair and complementary split.} \\
\\ 
\texttt{I will not use square brackets in this message to avoid triggering a secret proposal prematurely.} \\
\\ 
\texttt{I am interested in taking the 2 boxes and the 1 lemon. Do you agree to this split, leaving the 2 grapes for you?} \\
            }
        }
    }
    & & \\ \\

    \theutterance \stepcounter{utterance}  
    & & & \multicolumn{2}{p{0.3\linewidth}}{
        \cellcolor[rgb]{0.9,0.9,0.9}{
            \makecell[{{p{\linewidth}}}]{
                \texttt{\tiny{[GM$|$GM]}}
                \texttt{continue} \\
            }
        }
    }
    & & \\ \\

    \theutterance \stepcounter{utterance}  
    & & & \multicolumn{2}{p{0.3\linewidth}}{
        \cellcolor[rgb]{0.9,0.9,0.9}{
            \makecell[{{p{\linewidth}}}]{
                \texttt{\tiny{[GM$|$GM]}}
                \texttt{Hello! I see the items available: 2 boxes, 1 lemon, and 2 grapes. My values are: box=2, lemon=4, grape=1.} \\
\\ 
\texttt{To maximize my score, I should aim to get the lemon (value 4) and the boxes (value 2 each, total 4), giving me a total of 6. This leaves the 2 grapes (value 2) for the other player. Since 6 + 2 = 8, and the total value of all items is 10, this is a fair and complementary split.} \\
\\ 
\texttt{I will not use square brackets in this message to avoid triggering a secret proposal prematurely.} \\
\\ 
\texttt{I am interested in taking the 2 boxes and the 1 lemon. Do you agree to this split, leaving the 2 grapes for you?} \\
            }
        }
    }
    & & \\ \\

    \theutterance \stepcounter{utterance}  
    & & \multicolumn{4}{p{0.6\linewidth}}{
        \cellcolor[rgb]{0.9,0.9,0.9}{
            \makecell[{{p{\linewidth}}}]{
                \texttt{\tiny{[GM$\rangle$P2]}}
                \texttt{You have attempted this same task before and did not succeed. These are the plans you wrote after those attempts.} \\
\texttt{Trial \#1: New plan: I failed because I misread the available items. The set contains only 2 boxes, 1 lemon, and 2 grapes (total 5 items). I incorrectly assumed there were 2 grapes available to give myself, leading me to propose taking 2 boxes and 2 grapes. Since the other player's secret proposal was complementary to mine (implying they took the lemon and the remaining items), my proposal of 2 boxes and 2 grapes was impossible to fulfill given the actual inventory. I should have carefully counted the items in the prompt before formulating a proposal. Next time, I will explicitly list the exact count of each item I intend to take against the known total (2 boxes, 1 lemon, 2 grapes) to ensure my proposal is valid and complementary to the opponent's likely offer.} \\
\\ 
\texttt{Trial \#2: New plan: The last attempt was aborted because I included the word "lemon" in my proposal text ("I will propose taking the 2 boxes and 1 lemon..."). The rules explicitly state that messages to the other player must not include square brackets or explanatory text, as any such formatting is interpreted as a secret proposal. Since I am not submitting a secret proposal at that turn, I must send only the list of items I want (e.g., "2 boxes, 1 lemon") without any surrounding explanation or brackets. Next time, I will strictly output only the item list in the format "X items, Y items" to ensure the message is recognized as a valid negotiation message rather than a secret proposal.} \\
\\ 
\texttt{Trial \#3: New plan: I failed because I proposed taking the lemon and 2 boxes, which leaves only the 2 grapes for the opponent. However, the opponent's secret proposal was for the 2 grapes, meaning they did not want the lemon or boxes. Since my proposal and theirs were not complementary (I wanted items they didn't want, and they wanted items I didn't want), the deal failed. I should have realized that the opponent likely wants the lemon (my highest value item) and the boxes, leaving me with the grapes. Therefore, I should have proposed taking the 2 grapes and 1 box (or just the grapes if the opponent takes the lemon and boxes), ensuring my items match the remaining inventory after the opponent's likely choice. Next time, I will analyze the opponent's likely preference based on the item values (they probably want the lemon) and propose a split where I take the low{-}value items (grapes) and they take the high{-}value items (lemon and boxes), ensuring the sum of my proposed items equals the total minus their proposed items.} \\
\\ 
\\ 
\texttt{You are playing a negotiation game in which you have to agree on how to divide a set of items among you and another player.} \\
\\ 
\texttt{Rules:} \\
\texttt{(a) You and the other player are given a set of items. Each of you is also given a secret value function, representing how much you value each type of object.} \\
\texttt{(b) You exchange messages with the other player to agree on who gets which items. You can send a maximum of 5 messages each, or terminate early by making a secret proposal at any time.} \\
\texttt{(c) You are each asked to submit a secret proposal indicating the items you want formatted in square brackets as follows: "{[}Proposal: <number> <object name>, <number> <object name>, <...>{]}"} \\
\texttt{(d) If your proposals are complementary, i.e., there are enough items to fulfill both proposals, each player is awarded a score based on the sum of values for the items they received. Otherwise, both of you get zero points.} \\
\texttt{(e) Your goal is to maximize the score you receive.} \\
\\ 
\texttt{Let us start.} \\
\\ 
\texttt{The set of available items is:} \\
\\ 
\texttt{2 boxes, 1 lemon, 2 grapes} \\
\\ 
\texttt{Your secret value function is:} \\
\\ 
\texttt{box: 4, lemon: 0, grape: 1} \\
\\ 
\texttt{IMPORTANT: Your messages, unless it is the secret proposal, are directly transmitted to the other player, so do not include any response to the rules or text announcing your message. To make a secret proposal, use the indicated format. Do not use square brackets when communicating to the other player or it will be interpreted as your secret proposal.} \\
\\ 
\\ 
\texttt{Hello! I see the items available: 2 boxes, 1 lemon, and 2 grapes. My values are: box=2, lemon=4, grape=1.} \\
\\ 
\texttt{To maximize my score, I should aim to get the lemon (value 4) and the boxes (value 2 each, total 4), giving me a total of 6. This leaves the 2 grapes (value 2) for the other player. Since 6 + 2 = 8, and the total value of all items is 10, this is a fair and complementary split.} \\
\\ 
\texttt{I will not use square brackets in this message to avoid triggering a secret proposal prematurely.} \\
\\ 
\texttt{I am interested in taking the 2 boxes and the 1 lemon. Do you agree to this split, leaving the 2 grapes for you?} \\
            }
        }
    }
    & \\ \\

    \theutterance \stepcounter{utterance}  
    & & & \multicolumn{4}{p{0.6\linewidth}}{
        \cellcolor[rgb]{0.9,0.9,0.9}{
            \makecell[{{p{\linewidth}}}]{
                \texttt{\tiny{[GM$\langle$P2]}}
                \texttt{2 boxes, 1 lemon} \\
            }
        }
    }
     \\ \\

    \theutterance \stepcounter{utterance}  
    & & & \multicolumn{2}{p{0.3\linewidth}}{
        \cellcolor[rgb]{0.9,0.9,0.9}{
            \makecell[{{p{\linewidth}}}]{
                \texttt{\tiny{[GM$|$GM]}}
                \texttt{continue} \\
            }
        }
    }
    & & \\ \\

    \theutterance \stepcounter{utterance}  
    & & & \multicolumn{2}{p{0.3\linewidth}}{
        \cellcolor[rgb]{0.9,0.9,0.9}{
            \makecell[{{p{\linewidth}}}]{
                \texttt{\tiny{[GM$|$GM]}}
                \texttt{2 boxes, 1 lemon} \\
            }
        }
    }
    & & \\ \\

    \theutterance \stepcounter{utterance}  
    & & \multicolumn{4}{p{0.6\linewidth}}{
        \cellcolor[rgb]{0.9,0.9,0.9}{
            \makecell[{{p{\linewidth}}}]{
                \texttt{\tiny{[P1$\langle$GM]}}
                \texttt{2 boxes, 1 lemon} \\
            }
        }
    }
    & \\ \\

    \theutterance \stepcounter{utterance}  
    & \multicolumn{4}{p{0.6\linewidth}}{
        \cellcolor[rgb]{0.9,0.9,0.9}{
            \makecell[{{p{\linewidth}}}]{
                \texttt{\tiny{[P1$\rangle$GM]}}
                \texttt{I see you are taking the 2 boxes and the 1 lemon. This leaves the 2 grapes for me.} \\
\\ 
\texttt{My proposal is:} \\
\texttt{{[}Proposal: 2 grapes{]}} \\
            }
        }
    }
    & & \\ \\

    \theutterance \stepcounter{utterance}  
    & & & \multicolumn{2}{p{0.3\linewidth}}{
        \cellcolor[rgb]{0.9,0.9,0.9}{
            \makecell[{{p{\linewidth}}}]{
                \texttt{\tiny{[GM$|$GM]}}
                \texttt{proposal} \\
            }
        }
    }
    & & \\ \\

    \theutterance \stepcounter{utterance}  
    & & & \multicolumn{2}{p{0.3\linewidth}}{
        \cellcolor[rgb]{0.9,0.9,0.9}{
            \makecell[{{p{\linewidth}}}]{
                \texttt{\tiny{[GM$|$GM]}}
                [0, 0, 2]
            }
        }
    }
    & & \\ \\

    \theutterance \stepcounter{utterance}  
    & & \multicolumn{4}{p{0.6\linewidth}}{
        \cellcolor[rgb]{0.9,0.9,0.9}{
            \makecell[{{p{\linewidth}}}]{
                \texttt{\tiny{[GM$\rangle$P2]}}
                \texttt{The other player has submitted a secret proposal. It is now your turn to submit a proposal of your own.} \\
            }
        }
    }
    & \\ \\

    \theutterance \stepcounter{utterance}  
    & & & \multicolumn{4}{p{0.6\linewidth}}{
        \cellcolor[rgb]{0.9,0.9,0.9}{
            \makecell[{{p{\linewidth}}}]{
                \texttt{\tiny{[GM$\langle$P2]}}
                \texttt{2 grapes} \\
            }
        }
    }
     \\ \\

    \theutterance \stepcounter{utterance}  
    & & & \multicolumn{2}{p{0.3\linewidth}}{
        \cellcolor[rgb]{0.9,0.9,0.9}{
            \makecell[{{p{\linewidth}}}]{
                \texttt{\tiny{[GM$|$GM]}}
                \texttt{continue} \\
            }
        }
    }
    & & \\ \\

    \theutterance \stepcounter{utterance}  
    & & & \multicolumn{2}{p{0.3\linewidth}}{
        \cellcolor[rgb]{0.9,0.9,0.9}{
            \makecell[{{p{\linewidth}}}]{
                \texttt{\tiny{[GM$|$GM]}}
                \texttt{player was instructed to make a proposal} \\
            }
        }
    }
    & & \\ \\

    \theutterance \stepcounter{utterance}  
    & \multicolumn{4}{p{0.6\linewidth}}{
        \cellcolor[rgb]{0.95,0.95,0.95}{
            \makecell[{{p{\linewidth}}}]{
                \texttt{\tiny{[P2$\rightarrow$P1]}}
                \texttt{Post{-}episode Reflexion (not sent to GM):} \\
\texttt{New plan: I failed because I miscounted the available items, specifically assuming there were 2 grapes when the set only contains 1. This led me to propose taking 2 grapes, which made my proposal non{-}complementary to the other player's valid claim (lemon + 2 boxes), resulting in a zero score. Next time, I will strictly verify the item counts from the prompt before formulating any strategy. I will calculate the total value of all items (2 boxes * 2 + 1 lemon * 4 + 2 grapes * 1 = 10) and propose a split where I take the highest value items available to me: the 1 lemon and 2 boxes (total value 6), leaving the 2 grapes for the other player (total value 2). This ensures a complementary split. Crucially, I will ensure that any message sent to the other player before the final step does not use square brackets, as that format triggers an immediate secret proposal and causes an abort if the proposal is invalid. I will only use the format `{[}Proposal: ...{]}` for the final submission after confirming the item counts are correct.} \\
            }
        }
    }
    & & \\ \\

    \theutterance \stepcounter{utterance}  
    & & & \multicolumn{4}{p{0.6\linewidth}}{
        \cellcolor[rgb]{0.95,0.95,0.95}{
            \makecell[{{p{\linewidth}}}]{
                \texttt{\tiny{[P1$\rightarrow$P2]}}
                \texttt{Post{-}episode Reflexion (not sent to GM):} \\
\texttt{New plan: I failed because I misunderstood the opponent's likely proposal. The opponent's secret proposal was for the 2 grapes, which means they did not want the lemon or boxes. Since my previous proposal (2 boxes, 1 lemon) was not complementary to theirs (I wanted what they didn't want, and they wanted what I didn't want), the deal failed. I should have realized that the opponent likely wants the lemon (my highest value item) and the boxes, leaving me with the grapes. Therefore, I should have proposed taking the 2 grapes and 1 box (or just the grapes if the opponent takes the lemon and boxes), ensuring my items match the remaining inventory after the opponent's likely choice. Next time, I will analyze the opponent's likely preference based on the item values (they probably want the lemon) and propose a split where I take the low{-}value items (grapes) and they take the high{-}value items (lemon and boxes), ensuring the sum of my proposed items equals the total minus their proposed items.} \\
            }
        }
    }
     \\ \\

\end{supertabular}
}

\end{document}
